在平面界面 $z=0$ 两侧有两种无损各向同性非磁介质，入射侧 $z<0$ 折射率为 $n_1>n_2>0$，透射侧 $z>0$ 折射率为 $n_2$。真空波长为 $\lambda>0$ 的 TE（s）偏振单色平面波以入射角 $\theta$ 从介质1射向界面，假定 $\theta_c<\theta<\pi/2$，其中 $\theta_c=\arcsin(n_2/n_1)$。记 $k_0=2\pi/\lambda$。
(1) 从切向波矢守恒推导透射侧法向波矢成为虚数，并求倏逝场的 $1/e$ 振幅穿透深度；(2) 由 TE Fresnel 边界条件推导全内反射系数，证明模长为一并求反射相位；(3) 对窄角谱 TE 光束，用反射相位的驻相方法推导近似 Goos–Hänchen 横向位移；(4) 讨论临界角附近的穿透深度与位移、能流法向分量，以及近似公式何时不适用。
